% SPDX-License-Identifier: Apache-2.0
\section{\code{txhdl::pipeline} and \code{txhdl::funcs}}
\label{sec:r-operators}

Two modules that hold the same kind of thing, operators, and differ in
one property: whether the operator takes a cycle. An operator in
\code{pipeline} is \code{async} and awaits one edge of the process's
clock, so a pipeline's depth follows from the operators it awaits and
nobody places a stage boundary. An operator in \code{funcs} is a plain
\code{fn}, has no \code{.await} to write, and cannot count cycles, which
is the whole enforcement of the rule that a combinational body is
timeless. The operators with a Rust spelling are Rust's, on the value
types of Section~\ref{sec:r-types}; \code{funcs} holds the two that
have none, the arithmetic shift and the signed compare, and the
halves of a word.

\code{drive} is what makes an \code{async fn} a pipeline rather than
a sequence: at every edge at which an input is offered it starts an
invocation, it polls every invocation in flight as a process of its
own, each started as having crossed the edge it began at so its first
wait is the next edge, and it sends each result as its invocation
completes, oldest first. Several invocations are in flight at once,
each at a different await.

\lstinputlisting[caption={\code{lib/src/pipeline.rs}.},label={lst:r-pipeline}]{lib/src/pipeline.rs}
\lstinputlisting[caption={\code{lib/src/funcs.rs}.},label={lst:r-funcs}]{lib/src/funcs.rs}
